%!TEX root = ../main.tex
\chapter{Analiza bezpieczeństwa}
\label{ch:analiza-bezpieczenstwa}

\section{Model zagrożeń}
\label{sec:model-zagrozen}

\subsection{Przyjęte klasy adwersarzy}
\label{subsec:aktorzy}

Analiza opiera się na modelu adwersarza zdefiniowanym w sekcji~\ref{sec:zalozenia-adwersarz} i posługuje się wprowadzonymi tam pięcioma klasami, uporządkowanymi według rosnącego zakresu możliwości: A1 (adwersarz zewnętrzny), A2 (skompromitowany walidator dysponujący kluczem konsensusu nie więcej niż $f$ węzłów, gdzie $n = 3f + 1$), A3 (adwersarz sieciowy), A4 (adwersarz z kluczem domeny) oraz A5 (złośliwy operator węzła). Obowiązują również przyjęte tam granice modelu: adwersarz nie kontroluje więcej niż $f$ walidatorów, wykorzystywane prymitywy kryptograficzne pozostają bezpieczne, a infrastruktura DNS ofiary nie znajduje się pod jego kontrolą w chwili rejestracji. Konsekwencje naruszenia tych założeń omawia sekcja~\ref{sec:ograniczenia}.

\subsection{Aktywa i cele bezpieczeństwa}
\label{subsec:aktywa}

Analiza koncentruje się na ochronie czterech aktywów, których naruszenie podważyłoby użyteczność systemu jako warstwy zaufania. Stanowią one operacyjne ujęcie wymagań niefunkcjonalnych postawionych w sekcji~\ref{sec:wymagania-niefunkcjonalne}. To, co tam sformułowano jako cel projektowy, tu jest przedmiotem ochrony przed konkretnym adwersarzem:
\begin{enumerate}
    \item \textbf{integralność rejestru} - żaden podmiot nie może związać cudzej domeny z własnym kluczem ani zmienić bądź unieważnić cudzego wpisu;
    \item \textbf{niefałszowalność dowodów} - żaden dowód nie weryfikuje się względem uczciwego korzenia, jeśli nie odpowiada faktycznie zatwierdzonemu stanowi;
    \item \textbf{dostępność} - rejestracja, operacje modyfikujące oraz odczyty działają przy awarii co najwyżej $f$ walidatorów;
    \item \textbf{spójność replik} - wszyscy uczciwi walidatorzy utrzymują identyczny stan (skrót stanu) na tej samej wysokości łańcucha.
\end{enumerate}

\section{Klasyfikacja zagrożeń (STRIDE)}
\label{sec:stride}

Zagrożenia sklasyfikowano metodą STRIDE \cite{howard2003secure,shostack2014threat}, rozpatrując dla każdego z głównych komponentów systemu sześć kategorii: podszycie (\emph{Spoofing}), manipulację danymi (\emph{Tampering}), zaprzeczalność (\emph{Repudiation}), ujawnienie informacji (\emph{Information disclosure}), odmowę usługi (\emph{Denial of service}) oraz podniesienie uprawnień (\emph{Elevation of privilege}). Wynik klasyfikacji zestawiono w tabeli~\ref{tab:stride}.

\begin{table}[htbp]
    \centering
    \caption{Klasyfikacja zagrożeń STRIDE według komponentu (S - podszycie, T - manipulacja, R - zaprzeczalność, I - ujawnienie informacji, D - odmowa usługi, E - podniesienie uprawnień)}
    \label{tab:stride}
    \begin{tabular}{l c c c c c c}
        \toprule
        \textbf{Komponent}               & \textbf{S} & \textbf{T} & \textbf{R} & \textbf{I} & \textbf{D} & \textbf{E} \\
        \midrule
        Transakcja rejestracji           & \checkmark & \checkmark & -          & -          & \checkmark & \checkmark \\
        Transakcja rotacji/unieważnienia & \checkmark & \checkmark & -          & -          & \checkmark & \checkmark \\
        Stan (drzewo Merkle)             & -          & \checkmark & -          & -          & -          & -          \\
        Konsensus                        & \checkmark & \checkmark & -          & -          & \checkmark & \checkmark \\
        Weryfikacja domeny (DNS-01)      & \checkmark & \checkmark & -          & -          & \checkmark & -          \\
        Zapytania i dowody               & -          & \checkmark & -          & -          & \checkmark & -          \\
        \bottomrule
    \end{tabular}
\end{table}

Dwie kategorie pozostają w całej tabeli puste. \textbf{Zaprzeczalność} nie jest zagrożeniem z konstrukcji: każda operacja modyfikująca stan jest podpisana kluczem właściciela i trwale zapisana w łańcuchu, więc pełna, niezaprzeczalna historia operacji jest publicznie dostępna, a sam rejestr pełni rolę dziennika zdarzeń (sekcja~\ref{sec:koncepcja-dpki}). \textbf{Ujawnienie informacji} również nie jest zagrożeniem, ponieważ rejestr jest z założenia jawny, podobnie jak logi Certificate Transparency (sekcja~\ref{sec:ct}); nie zawiera danych poufnych poza kluczami prywatnymi, które nigdy nie opuszczają urządzenia właściciela i nie są częścią żadnej transakcji.

Oznaczone w tabeli zagrożenia posłużyły za punkt wyjścia do wyboru konkretnych wektorów ataku, które analizuje kolejna sekcja. Pojedynczy wektor odpowiada zwykle kilku komórkom tabeli naraz, ponieważ ten sam atak godzi jednocześnie w różne komponenty; rozważono ataki godzące w każde z czterech aktywów wskazanych w sekcji~\ref{subsec:aktywa}.

\section{Analiza wybranych wektorów ataku}
\label{sec:wektory-atakow}

Dla każdego rozważanego wektora przyjęto jednolity schemat: opis scenariusza ataku wraz ze wskazaniem klasy adwersarza, mechanizm obronny systemu oraz wynik testu empirycznego potwierdzającego ten mechanizm. Tam, gdzie atak nie był wykonalny w warunkach prototypu (np.\ ataki na poziomie konsensusu wymagające większej liczby walidatorów), analizę prowadzono na poziomie teoretycznym, wskazując właściwość gwarantującą obronę.

\subsection{Kompromitacja walidatora}
\label{subsec:kompromitacja-walidatora}

Najpoważniejszą klasą zagrożeń specyficznych dla rozproszonej architektury jest kompromitacja walidatora (adwersarze A2 i A5). Rozważono trzy scenariusze.

\paragraph{Manipulacja stanem z pominięciem aplikacji.} Adwersarz A5, mając dostęp do dysku własnego węzła, modyfikuje stan rejestru bezpośrednio w magazynie klucz--wartość, z pominięciem reguł aplikacji, podmieniając na przykład klucz publiczny przypisany do cudzej domeny. Chroni przed tym samo zobowiązanie kryptograficzne drzewa Merkle: jakakolwiek zmiana wartości liścia zmienia korzeń, a więc skrót stanu publikowany przez węzeł. Ma to dwie niezależne konsekwencje. Skrót stanu zmanipulowanego węzła przestaje zgadzać się ze skrótem uczciwej większości, przez co jego propozycje i głosy zostają odrzucone w konsensusie. Niezależnie od tego dowód wygenerowany ze zmanipulowanego drzewa nie zweryfikuje się względem uczciwego korzenia, a uczciwy dowód - względem korzenia zmanipulowanego. Manipulacja jest więc wykrywalna także po stronie klienta.

\paragraph{Kłamliwa odpowiedź węzła.} Skompromitowany węzeł (A2) udziela klientowi fałszywej odpowiedzi na zapytanie o stan domeny. Chroni przed tym model zaufania interfejsu klienta (sekcja~\ref{sec:api}): klient weryfikuje dołączony dowód względem korzenia pozyskanego z nagłówka bloku, niezależnie od tego, co zwrócił węzeł, a w trybie wielowęzłowym dodatkowo porównuje odpowiedzi i skróty stanu wielu węzłów na tej samej wysokości. Pojedynczy kłamliwy węzeł zostaje w ten sposób wykryty, o ile adwersarz nie kontroluje większości pozwalającej sfałszować również nagłówki, co wymaga ponad $f$ walidatorów i wykracza poza model.

\paragraph{Awaria walidatorów.} Scenariusze awarii przebadano empirycznie na referencyjnej sieci czterech walidatorów (sekcja~\ref{sec:siec-wielowezlowa}). Przy awarii jednego walidatora pozostała trójka dysponuje 75\% mocy głosów, a więc ponad wymaganymi dwiema trzecimi. Sieć kontynuowała produkcję bloków (wysokość łańcucha wzrosła w ciągu dziesięciu sekund obserwacji o cztery bloki, z~66 do~70), a rejestracja zlecona w tym czasie została zatwierdzona. Przy awarii dwóch walidatorów dostępne 50\% mocy głosów nie wystarcza do osiągnięcia kworum: konsensus wstrzymał się, czyli przez piętnaście sekund obserwacji nie powstał ani jeden nowy blok (wysokość pozostała na poziomie~71), zamiast dopuścić do rozwidlenia historii. Ilustruje to zasadę bezpieczeństwa ponad żywotnością (sekcja~\ref{subsec:bft-tendermint}): system woli przejściowo wstrzymać rejestracje, niż dopuścić do współistnienia dwóch sprzecznych widoków rejestru. Po ponownym uruchomieniu obu walidatorów odtworzyły one stan przez ponowne wykonanie bloków (sekcja~\ref{subsec:zarzadzanie-stanem}) i powróciły do konsensusu, osiągając skrót stanu identyczny z pozostałymi węzłami na wspólnej wysokości. Potwierdza to zachowanie spójności replik, czyli czwartego z aktywów wymienionych w sekcji~\ref{subsec:aktywa}.

\subsection{Ataki Sybil i Eclipse}
\label{subsec:sybil-eclipse}

Ataki na warstwę sieciową i tożsamość uczestników konsensusu wymagają odrębnego rozważenia, gdyż ich realizacja w warunkach prototypu była ograniczona.

\paragraph{Atak Sybil.} W ataku Sybil, opisanym po raz pierwszy przez Douceura \cite{douceur2002sybil}, adwersarz tworzy wiele fałszywych tożsamości, by uzyskać nieproporcjonalny wpływ na konsensus. W przyjętym modelu konsorcyjnym (\emph{permissioned}, sekcja~\ref{sec:rozproszony-rejestr}) atak ten jest z konstrukcji nieskuteczny: zbiór walidatorów jest jawny i ustalony w pliku konfiguracyjnym, a o sile głosu decyduje przynależność do tego zbioru, nie liczba tożsamości. Napastnik nie zwiększy więc swojego wpływu, powołując nowe węzły; jego możliwości pozostają ograniczone do $f$ walidatorów, które faktycznie kontroluje. Odróżnia to model konsorcyjny od sieci otwartych, w których odporność na Sybil wymaga kosztownego mechanizmu (dowód pracy lub dowód stawki).

\paragraph{Atak Eclipse i podział sieci.} Adwersarz sieciowy (A3), który izoluje część walidatorów, realizuje w istocie podział sieci. Żadna z powstałych części nie dysponuje ponad dwiema trzecimi mocy głosów, więc żadna nie sfinalizuje bloków samodzielnie. Skutek jest taki sam jak przy omówionej wyżej awarii dwóch walidatorów: sieć wstrzymuje produkcję bloków, zamiast rozdzielić się na dwie sprzeczne historie. Odrębnym przypadkiem jest atak Eclipse w postaci opisanej dla sieci Bitcoin przez Heilmana i innych \cite{heilman2015eclipse}, wymierzony w pojedynczego klienta, odizolowanego tak, by odpytywał wyłącznie węzły kontrolowane przez napastnika. Klient otrzyma spójne kłamstwo tylko wówczas, gdy adwersarz potrafi sfałszować również nagłówki bloków, co wymaga kontroli ponad $f$ walidatorów i leży poza modelem. Pełna obrona wymagałaby weryfikacji podpisów walidatorów nad nagłówkami przez w pełni funkcjonalnego lekkiego klienta, co wykracza poza zakres prototypu i wskazano jako kierunek rozwoju (sekcja~\ref{sec:ograniczenia}).

\paragraph{Atak długiego zasięgu.} Napastnik dysponujący historycznymi kluczami walidatorów (wariant A2) mógłby próbować odtworzyć alternatywną historię łańcucha od odległego punktu w przeszłości. Taki atak (\emph{long-range attack}), którego warianty i mechanizmy obrony zestawia przeglądowa praca Deirmentzoglou i innych \cite{deirmentzoglou2019longrange}, jest groźny w systemach opartych na dowodzie stawki, natomiast pozostaje bezprzedmiotowy przy zastosowanym konsensusie BFT o finalności deterministycznej (sekcja~\ref{subsec:pow-pos}): blok raz zatwierdzony jest ostateczny, a reorganizacja łańcucha nie jest dopuszczana, niezależnie od posiadanych kluczy historycznych.

\subsection{Odmowa usługi}
\label{subsec:odmowa-uslugi}

Odmowa usługi jest jedyną kategorią STRIDE osiągalną dla najsłabszego adwersarza (A1) bez żadnych dodatkowych uprawnień, dlatego sprawdzono ją empirycznie. Poza omówioną wyżej niedostępnością wynikającą z utraty kworum - czy to wskutek awarii walidatorów, czy podziału sieci - rozważono dwa warianty wyczerpywania zasobów węzła.

\paragraph{Zalew zniekształconych transakcji.} Aplikację poddano zalewowi 10\,000 zniekształconych transakcji o rozmiarze kilobajta każda. Wszystkie odrzucono na etapie dekodowania, jeszcze przed kosztownymi operacjami kryptograficznymi, w łącznym czasie 32~ms (średnio 3,2~µs na transakcję). Rejestracja zlecona bezpośrednio po zalewie przeszła bez zakłóceń, a więc koszt odrzucenia śmieciowego ruchu nie przełożył się na dostępność operacji dla uczciwych użytkowników.

\paragraph{Transakcje z nadmiernie długimi polami.} Transakcje z polami rzędu megabajta (nazwa domeny, powód unieważnienia) są odrzucane przez walidację składniową, zanim dotkną stanu, dzięki limitom długości: 253 bajty dla nazwy domeny i 256 bajtów dla powodu unieważnienia. Koszt takiego ataku po stronie węzła ogranicza się więc do samej walidacji składniowej.

\subsection{Próba podszycia się pod właściciela domeny}
\label{subsec:podszycie-domeny}

Najbliżej istoty PKI leży próba związania cudzej domeny z kluczem napastnika lub przejęcia istniejącego wiązania. Rozważono kilka wariantów, wszystkie w zasięgu najsłabszego adwersarza A1.

\paragraph{Wyprzedzenie rejestracji (\emph{front-running}).} Adwersarz obserwuje oczekującą transakcję rejestracji i opublikowany dla niej rekord TXT z wartością wyzwania, po czym próbuje zarejestrować tę samą domenę na własny klucz, wykorzystując zaobserwowaną wartość. Chroni przed tym związanie wartości wyzwania z kluczem publicznym (sekcja~\ref{sec:wlasnosc-domeny}). Węzeł wylicza oczekiwaną wartość z klucza nadawcy transakcji, więc dla napastnika wypada ona inaczej niż wartość opublikowana przez prawdziwego właściciela. Opublikowany rekord TXT nie pasuje więc do jego transakcji i rejestracja adwersarza zostaje odrzucona.

\paragraph{Nieautoryzowana modyfikacja cudzego wpisu.} Napastnik próbuje unieważnić lub przejąć cudzą domenę, składając transakcję rotacji bądź unieważnienia podpisaną własnym kluczem. Zapobiega temu reguła autoryzacji (sekcja~\ref{sec:protokoly}): węzeł odczytuje klucz uprawniający do modyfikacji ze stanu rejestru. Wartość podana w transakcji jest ignorowana, więc podpis musi pochodzić od bieżącego właściciela.

\paragraph{Przekierowanie podpisanej transakcji.} Napastnik przechwytuje poprawnie podpisaną transakcję unieważnienia jednej domeny i podmienia w niej nazwę domeny na inną, licząc na to, że podpis pozostanie ważny. Podpis obejmuje jednak całą treść transakcji, w tym nazwę domeny (sekcja~\ref{subsec:reprezentacja-cert}), więc jej zmiana zmienia bajty do podpisania i unieważnia podpis. Tak zmodyfikowana transakcja zostaje odrzucona.

\paragraph{Ponowna rejestracja istniejącej domeny.} Napastnik próbuje nadpisać istniejące wiązanie, rejestrując zajętą już domenę na własny klucz. Zapobiega temu warunek nieistnienia domeny sprawdzany przy rejestracji: próba zostaje odrzucona z kodem błędu semantycznego, a klucz dotychczasowego właściciela pozostaje w stanie nienaruszony.

\subsection{Ataki na warstwę kryptograficzną i protokół}
\label{subsec:ataki-kryptograficzne}

Na końcu rozważono ataki wymierzone w warstwę kryptograficzną i w protokół przetwarzania transakcji.

\paragraph{Fałszowanie podpisu.} Próby przedłożenia transakcji z podpisem sfałszowanym, zmodyfikowanym bądź w ogóle brakującym, a także transakcji, w której klucz w certyfikacie nie odpowiada kluczowi podpisującemu, sprowadzają się do złamania własności schematu Ed25519 nad deterministycznymi bajtami do podpisania (sekcja~\ref{subsec:reprezentacja-cert}). Przy założeniu bezpieczeństwa prymitywu są one nieskuteczne, co potwierdza zestaw testów obejmujący sfałszowany i brakujący podpis oraz niezgodność klucza certyfikatu z kluczem podpisującym.

\paragraph{Ataki powtórzeniowe.} Atak powtórzeniowy (\emph{replay attack}) polega na ponownym przedłożeniu przechwyconej, poprawnie podpisanej transakcji. Rozważono trzy jego warianty.

Pierwszym jest powtórzenie transakcji modyfikującej w tym samym łańcuchu. Zapobiega mu ściśle rosnący licznik operacji (\emph{nonce}). Po wykonaniu transakcji jej licznik traci ważność, więc te same bajty nie zostaną przyjęte po raz drugi. Drugim wariantem jest powtórzenie transakcji w innym łańcuchu. Zapobiega mu identyfikator łańcucha wchodzący w skład podpisywanych bajtów, dzięki któremu podpis złożony dla jednego łańcucha nie jest ważny w żadnym innym. Trzeci wariant polega na powtórzeniu pierwotnej transakcji rejestracji już po unieważnieniu domeny, aby przywrócić unieważnione wiązanie; scenariusz ten jest szczególnie istotny przy kompromitacji klucza właściciela (adwersarz A4). Zapobiega mu trwałość unieważnienia (sekcja~\ref{subsec:uniewaznienie}): unieważniony wpis nie znika ze stanu, więc domena wciąż istnieje i nie spełnia warunku nieistnienia wymaganego przy rejestracji. Powtórzona transakcja zostaje zatem odrzucona, mimo że jej podpis jest poprawny.

Spośród dwudziestu dwóch rozważonych wektorów dziewiętnaście potwierdzono empirycznie: testami bezpieczeństwa w pakiecie \texttt{internal/app} oraz scenariuszami uruchamianymi na sieci czterech walidatorów skryptem \texttt{scripts/test-testnet.sh}. Pozostałe trzy - atak Sybil, partycję sieci wraz z atakiem Eclipse oraz atak długiego zasięgu - przeanalizowano teoretycznie, ponieważ wymagają kontroli nad warstwą sieciową bądź zbiorem walidatorów w skali niedostępnej w prototypie. Nazwy odpowiadających im funkcji testowych i scenariuszy, wraz z możliwością ich powtórzenia, udostępnia repozytorium projektu wskazane w sekcji~\ref{sec:testy}.

\section{Porównanie z klasyczną PKI}
\label{sec:porownanie-bezpieczenstwa}

Zestawienie właściwości bezpieczeństwa zaproponowanego systemu z klasyczną PKI opartą na X.509 i urzędach certyfikacji przedstawia tabela~\ref{tab:porownanie-bezp}. Porównanie odnosi się do słabości strukturalnych zidentyfikowanych w rozdziale~\ref{ch:pki-klasyczna}.

\begin{table}[htbp]
    \centering
    \caption{Porównanie właściwości bezpieczeństwa klasycznej PKI i zaproponowanego systemu DPKI}
    \label{tab:porownanie-bezp}
    \renewcommand{\arraystretch}{1.15}
    \begin{tabular}{L{0.27\textwidth} L{0.31\textwidth} L{0.31\textwidth}}
        \toprule
        \textbf{Właściwość}                  & \textbf{Klasyczna PKI (X.509 + CA)}                 & \textbf{DPKI}                                                      \\
        \midrule
        Punkt zaufania                       & pojedyncze CA - każde z setek w magazynie zaufania  & ponad dwie trzecie mocy głosów walidatorów                         \\
        \addlinespace
        Skutek kompromitacji punktu zaufania & fałszywy certyfikat dla dowolnej domeny (DigiNotar) & poniżej $1/3$: brak skutku; do $2/3$: wstrzymanie, nie fałszerstwo \\
        \addlinespace
        Unieważnienie                        & CRL/OCSP - opóźnione, w praktyce \emph{soft-fail}   & następny blok, trwałe, potwierdzane dowodem                        \\
        \addlinespace
        Dowód statusu                        & łańcuch certyfikatów oraz odpowiedź OCSP            & dowód Merkle względem skrótu stanu, weryfikowalny offline          \\
        \addlinespace
        Dowód nieistnienia                   & brak                                                & dowód nieobecności                                                 \\
        \addlinespace
        Transparentność                      & Certificate Transparency jako opcjonalny dodatek    & wbudowana: rejestr jest dziennikiem                                \\
        \addlinespace
        Wykrywalność nadużycia               & monitoring logów CT, po fakcie                      & nadużycie nie zostaje zatwierdzone (prewencja)                     \\
        \bottomrule
    \end{tabular}
\end{table}

Zasadnicza różnica dotyczy skutku kompromitacji punktu zaufania. W modelu klasycznym kompromitacja dowolnego z setek urzędów pozwala wystawić ważny certyfikat dla dowolnej domeny; o bezpieczeństwie decyduje najsłabsze ogniwo (sekcja~\ref{subsec:spof}). W modelu zaproponowanym kompromitacja mniejszości walidatorów nie wywołuje żadnego skutku dla integralności rejestru, a osiągnięcie progu dwóch trzecich prowadzi co najwyżej do wstrzymania systemu, nie do wystawienia fałszywego wiązania. Druga różnica dotyczy unieważnień: scentralizowane, opóźnione i podatne na \emph{soft-fail} mechanizmy CRL/OCSP (sekcja~\ref{sec:crl-ocsp}) zastąpiono unieważnieniem propagującym się w czasie jednego bloku, trwałym i potwierdzanym dowodem względem skrótu stanu. Transparentność, którą Certificate Transparency dodaje do klasycznej PKI jako warstwę wykrywania po fakcie (sekcja~\ref{sec:ct}), jest tu wbudowana w samą konstrukcję rejestru i ma charakter prewencyjny: nieautoryzowane wiązanie nie zostaje zatwierdzone, a nie jedynie wykryte po wystawieniu.

\section{Ograniczenia zaproponowanego rozwiązania}
\label{sec:ograniczenia}

Trzeba też wskazać zagrożenia, przed którymi system \emph{nie} chroni, oraz założenia, których naruszenie podważa jego gwarancje. Wyróżniono cztery takie ograniczenia.

\paragraph{Cenzura przez większość walidatorów.} Koalicja kontrolująca ponad dwie trzecie mocy głosów może odmawiać włączania wybranych transakcji do bloków, skutecznie blokując rejestrację lub modyfikację konkretnej domeny. Zagrożenie to jest niewykrywalne kryptograficznie i wynika z samej konstrukcji konsensusu BFT; ograniczyć je można jedynie na poziomie zarządzania, przez dobór zróżnicowanego, niezależnego zbioru walidatorów.

\paragraph{Kompromitacja DNS w oknie rejestracji.} Mechanizm DNS-01 dowodzi kontroli nad strefą DNS domeny, a nie ,,własności'' domeny w sensie prawnym. Adwersarz kontrolujący strefę DNS ofiary w momencie rejestracji może zarejestrować domenę na własny klucz. Ograniczenie to jest wspólne z całym ekosystemem ACME i wynika z zakorzenienia tożsamości w istniejącym DNS; jego usunięcie wymagałoby innego źródła tożsamości domeny.

\paragraph{Utrata klucza domeny przed rotacją.} Adwersarz, który wszedł w posiadanie klucza prywatnego właściciela (A4), może unieważnić wiązanie lub przejąć je rotacją, zanim prawowity właściciel zareaguje. System gwarantuje w tym przypadku jedynie, że każda taka operacja jest publicznie i trwale odnotowana (wykrywalność), oraz że unieważnienia nie da się cofnąć; nie zapobiega jednak samemu nadużyciu kluczem, który z definicji jest jedynym uprawnieniem do operacji.

\paragraph{Brak pełnego lekkiego klienta.} Tryby weryfikacji udostępnione klientowi (sekcja~\ref{sec:api}) weryfikują dowód względem korzenia pozyskanego z nagłówka bloku, lecz w obecnej implementacji nie weryfikują podpisów walidatorów nad samym nagłówkiem. Pełna minimalizacja zaufania w scenariuszu ataku Eclipse na klienta (sekcja~\ref{subsec:sybil-eclipse}) wymaga w pełni funkcjonalnego lekkiego klienta sprawdzającego te podpisy; jego implementację wskazano jako kierunek dalszych prac (rozdział~\ref{ch:podsumowanie}).

\paragraph{} Analiza wykazała, że system usuwa najpoważniejszą słabość klasycznej PKI - kompromitację pojedynczego punktu zaufania - i odpiera badane ataki na warstwę kryptograficzną, protokół oraz konsensus. Cenę stanowią dwa nowe założenia: próg dwóch trzecich uczciwych walidatorów oraz zaufanie do DNS jako źródła tożsamości. Wydajność, koszt operacji i skalowalność systemu omawia następny rozdział.
